%! Zeros of Polynomials: the Rational Root Theorem — Unit 4, Lesson 4.4 — Homework
\documentclass[10pt]{article}
\usepackage{algebra2-article}
\usepackage{algebra2-boxes}
\newcommand{\TallMath}[1]{$\displaystyle #1\rule[-1.4em]{0pt}{3.2em}$}

\begin{document}

\pageheader{Unit 4, Lesson 4.4}{Homework}
\namedateperiod

\vspace{0.10in}

\begin{notesbox}{Practice}
Constant term on top, leading coefficient on the bottom, and every candidate gets a $\pm$. When you
find a zero, divide it out and \emph{finish} --- an answer is not done until every factor is prime.

\begin{enumerate}[leftmargin=1.9em, itemsep=11pt]
  \item \textbf{Just the list.} Write every possible rational zero. Reduce and remove repeats.
    \begin{enumerate}[label=(\alph*), leftmargin=1.8em, itemsep=5pt]
      \item $f(x) = x^3 + 3x^2 - 6x - 8$: \blank{6.4cm}
      \item $g(x) = 2x^3 - 5x^2 - 4x + 3$: \blank{6.4cm}
      \item $h(x) = 6x^3 + x^2 - 5x - 2$: \blank{6.4cm}

        \vspace{2pt}
        How many candidates are on list (c)? \blank{1.2cm}
    \end{enumerate}

  \item \textbf{Leading coefficient $1$.} Find all the zeros. Show the synthetic division on your own
    paper.
    \begin{enumerate}[label=(\alph*), leftmargin=1.8em, itemsep=5pt]
      \item $f(x) = x^3 + 3x^2 - 6x - 8$: \; first zero $r = $ \blank{1.2cm}, depressed polynomial
        \blank{2.8cm}

        \vspace{2pt}
        Completely factored: \blank{4.6cm} \quad Zeros: \blank{3.2cm}
      \item $p(x) = x^3 - 6x^2 + 3x + 10$: \; completely factored \blank{4.6cm}

        \vspace{2pt}
        Zeros: \blank{3.2cm}
    \end{enumerate}

  \item \textbf{Leading coefficient not $1$.} Solve $2x^3 - 5x^2 - 4x + 3 = 0$.
    \begin{enumerate}[label=(\alph*), leftmargin=1.8em, itemsep=5pt]
      \item A candidate that works: $r = $ \blank{1.2cm} \; (show your test: \blank{3.4cm})
      \item Depressed polynomial: \blank{2.8cm} \; which factors as \blank{2.8cm}
      \item Completely factored: \blank{5.0cm}
      \item Solutions: \blank{3.6cm} \; Which one would an integers-only list have missed?
        \blank{1.4cm}
    \end{enumerate}

  \item \textbf{Degree 4.} Solve $x^4 - 4x^3 - x^2 + 16x - 12 = 0$. You will need the workflow
    \emph{twice}.
    \begin{enumerate}[label=(\alph*), leftmargin=1.8em, itemsep=5pt]
      \item Candidates: \blank{5.4cm}
      \item First zero $r_1 = $ \blank{1.0cm}, depressed polynomial \blank{3.0cm}
      \item Second zero $r_2 = $ \blank{1.0cm}, next depressed polynomial \blank{2.8cm}
      \item Completely factored: \blank{5.0cm} \quad Solutions: \blank{3.6cm}
    \end{enumerate}

  \item \textbf{No hints this time.} On Lesson 4.3's exit ticket you were \emph{told} to divide
    $x^3 + 2x^2 - 5x - 6$ by $(x+1)$. Today nobody tells you.
    \begin{enumerate}[label=(\alph*), leftmargin=1.8em, itemsep=5pt]
      \item Candidates: \blank{5.4cm}
      \item Solve $x^3 + 2x^2 - 5x - 6 = 0$: \blank{4.0cm}
      \item \textbf{Verify} your answer two ways: substitute one solution to show it gives $0$
        \blank{4.0cm}, and state what the graph of this function must look like at each solution
        \blank{4.6cm}
    \end{enumerate}

  \item \textbf{Not every zero is rational.} Solve $x^3 - 2x^2 - 3x + 6 = 0$.
    \begin{enumerate}[label=(\alph*), leftmargin=1.8em, itemsep=5pt]
      \item Candidates: \blank{4.6cm} \quad The one that works: $r = $ \blank{1.0cm}
      \item Depressed polynomial: \blank{2.4cm} \quad Solve it: $x = $ \blank{2.4cm}
      \item Two of the three solutions were never on your candidate list. Explain why the Rational
        Root Theorem could not have found them. \blank{6.0cm}
    \end{enumerate}

  \item \textbf{Explain.} Answer each in one or two sentences.
    \begin{enumerate}[label=(\alph*), leftmargin=1.8em, itemsep=5pt]
      \item Why is the candidate list for a polynomial with leading coefficient $1$ made only of
        integers? \writelines{2}
      \item A classmate tests every candidate on a cubic and none of them gives $0$. They conclude
        that they made an arithmetic mistake. Are they necessarily right? \writelines{2}
    \end{enumerate}
\end{enumerate}
\end{notesbox}

\begin{extensionbox}
\textbf{Build them, don't just break them.}
\begin{enumerate}[label=(\alph*), leftmargin=1.8em, itemsep=5pt]
  \item A rectangular planter box holds $V(x) = 2x^3 + 3x^2 - 3x - 2$ cubic feet, where $x$ is
        measured in feet. Find all three dimensions as linear factors, then explain why $x$ must be
        greater than $1$ for the box to be physically possible.
  \item Check your answer to (a) at $x = 3$: compute $V(3)$ directly, then multiply your three
        dimensions at $x = 3$. Do they agree?
  \item Write a cubic with \emph{integer} coefficients whose zeros are $-\tfrac23$, $1$, and $4$.
        (Start from the factors $(3x+2)$, $(x-1)$, and $(x-4)$.) Then confirm that $-\tfrac23$ really
        does appear on your polynomial's own candidate list.
  \item A polynomial with integer coefficients has $\tfrac{3}{5}$ as a zero. What is the smallest
        possible value its leading coefficient could have, and why?
\end{enumerate}
\writelines{4}
\end{extensionbox}

\begin{spiralbox}
\textbf{Coming next --- Lesson 4.5: The Fundamental Theorem of Algebra \& Complex Zeros.} Twice today
a cubic gave up only one or two real solutions and then stalled: $x^3-2x^2-3x+6$ needed
$\pm\sqrt{3}$ to reach three, and a depressed polynomial like $x^2+4$ has no real solutions at all.
Tomorrow you find out that this is not a shortage --- a degree-$n$ polynomial has \emph{exactly} $n$
solutions once imaginary numbers are allowed, they arrive in conjugate pairs, and you can run the
whole process backward to build a polynomial from a list of zeros.
\end{spiralbox}

\end{document}
